\section{Accessibility claims}
\label{sec:access}

DiGiTerra's guided page is built to WCAG 2.2 AA: landmarks, a label for every control, live regions for errors and progress, focus that moves only where the user's own action opened the next step, and no meaning carried by color alone. A Chromium test runs axe-core under the WCAG A and AA rule sets with no serious violation allowed, reads a typed hyperparameter back from the fitted estimator, and walks the page by keyboard; \code{docs/ACCESSIBILITY.md} states what is checked and what is not claimed, and \code{docs/guided\_page\_aria\_transcript.txt} is the accessible-name transcript of the page. What is not claimed is a screen-reader session with a human user; the checks are automated.

Among the compared products, the feature survey found three kinds of statement. The cloud vendors publish conformance reports: AWS an ACR/VPAT covering Section 508 and WCAG for SageMaker, Microsoft ACRs for WCAG, Section 508 and EN 301 549, and Google a VPAT list that names Vertex AI Tabular Data. Two of the free desktop tools publish their own statements: jamovi's says the interface ``can be operated entirely from the keyboard and works with NVDA and JAWS'', and JASP's VPAT of August 2025 says ``JASP does not currently support screen readers''. Orange, KNIME, Weka, AI Studio and DataRobot state nothing on the pages read; H2O Flow documents keyboard shortcuts only. A statement is not a test, and this section records the statements.
